\documentclass[10pt,a4paper]{amsart}
\usepackage[margin=19mm]{geometry}
\usepackage{amsmath,amssymb,mathtools}
\usepackage[scr=rsfs,cal=euler]{mathalfa}
\usepackage{xfrac}
\usepackage[unicode,hidelinks,bookmarks=false]{hyperref}
\newcommand{\ie}{i.e.\ }
\newcommand{\cf}{cf.\ }
\newcommand{\resp}{respectively\ }
\newcommand{\Subsection}{\subsection}
\providecommand{\nobreakdash}{\nobreak\hbox{-}}
\setlength{\emergencystretch}{2em}
\multlinegap0pt
\usepackage[shortlabels]{enumitem}
\usepackage{amssymb}
\usepackage{xcolor}
\usepackage{booktabs}
\usepackage{caption}
\usepackage{float}
\usepackage{natbib}
\newtheorem{lem}{Lemma}
\newcommand{\E}{\mathbb{E}}
    \newcommand{\M}{\mathcal{M}}
\newcommand{\bX}{\mathbf{X}}
\newcommand{\bbeta}{\boldsymbol{\beta}}
\providecommand{\tightlist}{%
  \setlength{\itemsep}{0pt}\setlength{\parskip}{0pt}}\usepackage{longtable,booktabs,array}
\title{\bf Improving variable selection properties with data integration and transfer learning}
\author{Paul Rognon-Vael, David Rossell, Piotr Zwiernik}
\date{Source: arXiv:2502.15584v3 (CC BY 4.0)}
\begin{document}
\maketitle

\noindent\emph{Source and licence.} \bf Improving variable selection properties with data integration and transfer learning. Version arXiv:2502.15584v3 (CC BY 4.0), distributed by arXiv under the Creative Commons Attribution 4.0 International licence, http://creativecommons.org/licenses/by/4.0/, as recorded in the arXiv licence metadata for this exact version and shown on the abstract page https://arxiv.org/abs/2502.15584v3. Full licence notice: \texttt{licenses/arxiv-2502.15584-CC-BY-4.0.txt}.\par\smallskip

\noindent\textit{Editorial scope.} Counted unit: the article's lem lem:boundexpnc (source0.tex lines 1292--1294). The article's complete original proof of it is delivered with it (source0.tex lines 1296--1349, one \texttt{proof} environment of 54 lines). No mathematics is written, repaired or re-derived: the statement and the proof are the article's own lines, in the article's own notation, and no retained line is substituted or re-flowed.  Retained uncounted as the source-local context the proof needs: lines 410--417; lines 1117--1120; lines 1156--1160; lines 1211--1220; lines 1216--1218; lines 1280--1282.  Nothing was omitted from any retained span, so the package carries no omission record.  No retained line contains a citation, so the wrapper carries no bibliography and no bibliography entry.  No retained line contains a figure environment, an image-inclusion command or drawing code, so no image is delivered and the package declares no asset dependency.  Editorial formatting only, in addition: an AMS \texttt{amsart} preamble that restates the article's own declarations and macros used by the retained lines, namely the environments \texttt{lem} on separate counters, together with the standard packages those lines need.  The retained environments print in this document as Lemma 1 in order of appearance, whereas the article prints the same items with the numbering of its own full document.  The complete native source of the article as distributed in the arXiv e-print for this exact version is bundled unchanged as \texttt{sources/173658/source0.tex}.
\begin{lem}\label{lem:L0toL1convergence}
\begin{enumerate}[label= (\roman*)]\tightlist
\item
$P(\hat{S}^{ei} \neq S) \,\,\leq \,\, 2 \sum_{M \in \M\setminus \{S\}} \E\left(NC(M)\right)$.
\item  For any $M,M'\in\M$, such that $M\neq M'$, 
$NC(M)\;\leq \;\big(1+\exp(C(M')-C(M))\big)^{-1}$.
\end{enumerate}
\end{lem}

\begin{lem}\label{lemma:s2} Let $W \sim \chi_\nu^2(\mu)$ with $\mu>0$. For any $w<\mu$, 
$$
P(W<w) \leq \frac{e^{-\frac{1}{2}(\sqrt{\mu}-\sqrt{w})^2}}{(\mu / w)^{\nu / 4}}.$$
\end{lem}

\begin{lem}\label{lemma:s20}
 Let $W\sim \chi_\eta^2(\mu)$ with $\mu \geq 0$. Assume that $h$, $c$, $\eta$ and $\mu$ are functions of $n$ such that $h$ is positive and increasing, $c$ is positive such that $\lim_{n\to\infty}c = l \in (0,\infty)$, $\eta=o(\ln(h))$, and $\mu=o(\ln(h))$. Then for any $\alpha\in(0,\min\{1,l/2\})$ and every large enough $n$, we have
$$
\int_0^1 P\left(W>c \ln \left(\frac{h}{1 / u-1}\right)\right) d u = o\big(  h^{-\alpha}\big).$$
\end{lem}

\begin{lem}\label{lemma:s7}
    Let $M,Q$ be any two nested models such that $M \subseteq Q$. Then
$$
L_{QM}\;=\;\|\bX_{Q}\tilde{\bbeta}^{(Q)}\|^2-\|\bX_{M}\tilde{\bbeta}^{(M)}\|^2\;\sim\; \chi_{|Q\setminus M|}^2\left(\mu_{Q M}\right), 
$$
where $\chi_{k}^2(\mu)$ denotes, when $\mu>0$, the noncentral chi-squared distribution with $k$ degrees of freedom and noncentrality parameter $\mu$ and, when $\mu=0$, the chi-squared distribution with $k$ degrees of freedom $\chi_{k}^2$. The parameter $\mu_{QM}$ is given by \begin{equation}\label{eq:muQM}
    \mu_{QM}\;\;:=\;\;\|\big(I_n-P_M\big) \bX_{Q\setminus M} \bbeta^*_{Q\setminus M}\|^2
\end{equation}
where $P_M=\bX_{M}\left(\bX_{M}^{\top} \bX_{M}\right)^{-1} \bX_{M}^{\top}$.
\end{lem}

where $\chi_{k}^2(\mu)$ denotes, when $\mu>0$, the noncentral chi-squared distribution with $k$ degrees of freedom and noncentrality parameter $\mu$ and, when $\mu=0$, the chi-squared distribution with $k$ degrees of freedom $\chi_{k}^2$. The parameter $\mu_{QM}$ is given by \begin{equation}
    \mu_{QM}\;\;:=\;\;\|\big(I_n-P_M\big) \bX_{Q\setminus M} \bbeta^*_{Q\setminus M}\|^2
\end{equation}

\begin{equation}\label{eq:DeltaM}
\Delta_{MT}\;:=\;\sum_{j=1}^b \kappa_j(|M_j|-|T_j|) \quad\text{and} \quad L_{TM}\;:=\;\|\bX_{T}\tilde{\bbeta}^{(T)}\|^2-\|\bX_{M}\tilde{\bbeta}^{(M)}\|^2.
\end{equation}

\begin{lem}\label{lem:boundexpnc} For any $T\subseteq S$ and $M\in \M\setminus\{T\}$, denote $Q_T= M\cup T$ and $A_T:=\gamma\Delta_{MT}+\tfrac{1-\gamma}{6}\mu_{Q_TM}$ (\emph{cf}~\eqref{eq:DeltaM} and~\eqref{eq:muQM}). Suppose that, for some $\gamma \in (1/2,1]$, it holds that $A_T>0$, $|M\setminus T|=o(A_T)$, and $\mu_{Q_T T}=o(A_T)$. For any $\psi\in(0,1)$ and every $n$ large enough,
$$\E\left(NC(M)\right)  \;\;\leq\;\; e^{-\psi A_T}.$$
\end{lem}

\begin{proof}
Since $M\in \mathcal M\setminus \{T\}$, by Lemma~\ref{lem:L0toL1convergence} (ii),  $NC(M)\leq (1+e^{C(T)-C(M)})^{-1} \in [0,1]$.  
The first step of the proof is to use that for any random variable $Z \geq 0$ we have $\E(Z)=\int_0^\infty P(Z>u) du$, so that
\begin{eqnarray*}
  \E\left(NC(M)\right)  &\leq& \int_0^1 P\left((1+e^{C(T)-C(M)})^{-1}\geq u\right){\rm d} u\\
  &=&\int_0^1 P\left(C(T)-C(M) \leq \ln{(\tfrac1u-1)}\right){\rm d}u\\
  &=&\int_0^1 P\Big(-\tfrac12 L_{TM} \geq \Delta_{MT}-\ln{(\tfrac1u-1)}\Big) {\rm d}u.
\end{eqnarray*}

The second step of the proof is to use the union bound to upper bound the probability in the integrand above.
Let $Q_T=M \cup T$ and recall that $L_{TM}=L_{Q_TM}-L_{Q_TT}$. For any $\gamma$ 
$$
-\tfrac12 L_{TM} -(\Delta_{MT}-\ln{(\tfrac1u-1)})\;=\;\left(\tfrac12 L_{Q_TT}-\ln{\Bigl(\frac{e^{\gamma\Delta_{MT}}}{\tfrac1u-1}\Bigl)}\right)-\Bigl(\tfrac12 L_{Q_TM}+(1-\gamma)\Delta_{MT}\Bigl).
$$
Observe that for any random variables $U$, $V$, and any $\epsilon,\gamma' \geq 0$, the event $
\{U-V \geq 0\}$ implies $\{U \geq \gamma' \epsilon\}\cup
\{V  < \gamma' \epsilon\}$. Let $U=\tfrac12 L_{Q_TT}-\ln{\Bigl(\frac{e^{\gamma\Delta_{MT}}}{1/u-1}\Bigl)}$ and $V=\tfrac12 L_{Q_TM}+(1-\gamma)\Delta_{MT}$. Take $\epsilon=\mu_{Q_TM}$ and $\gamma'=\tfrac16({1-\gamma})$, and observe that
$A_T:=\gamma\Delta_{MT}+\tfrac{1-\gamma}{6}\mu_{Q_TM}= \gamma\Delta_{MT}+\gamma'\mu_{Q_TM}$. We then have
\begin{align*}
\{U\geq \gamma'\epsilon\}\;&=\;\{\tfrac12 L_{Q_TT}\geq \ln{\Bigl(\frac{e^{\gamma\Delta_{MT}}}{\tfrac1u-1}}\Bigl)+\gamma'\mu_{Q_TM}\}\;=\;\{\tfrac12 L_{Q_TT}\geq \ln{\Bigl(\frac{e^{A_T}}{\tfrac1u-1}}\Bigl)\}\\
\{V<\gamma'\epsilon\}\;&=\;\{\tfrac12 L_{Q_TM}<-(1-\gamma)\Delta_{MT}+\gamma'\mu_{Q_TM}\}\;=\;\{\tfrac12 L_{Q_TM}<\gamma'(\mu_{Q_TM}-6\Delta_{MT})\}.
\end{align*}
By the union bound we have that
\begin{equation}\label{eq:prooflemmas20nonspu}
    \E\left(NC(M)\right)  \leq \int_0^1P\Bigl(\tfrac12 L_{Q_TT}\geq \ln{\Bigl(\frac{e^{A_T}}{\tfrac1u-1}\Bigl)}\Bigl){\rm d}u+ P\Bigl(\tfrac12 L_{Q_TM}<\gamma'(\mu_{Q_TM}-6\Delta_{MT})\Bigl).
\end{equation}

The third and final step of the proof is to upper bound each of the terms in the right-hand side of~\eqref{eq:prooflemmas20nonspu}. The intuition is that both $T$ and $M$ are nested within $Q_T$, and therefore $L_{Q_T T}$ and $L_{Q_T M}$ follow chi-squared distributions.
We first bound the first term. 
If $M \subset T$ then $Q_T=T$, $L_{Q_TT}=0$, and this term is zero. Suppose now that $M \not\subset T$.  
Then, by Lemma~\ref{lemma:s7}, 
$L_{Q_TT} \sim \mathcal{\chi}_{|Q_T\setminus T|}^2(\mu_{Q_TT})$ with $|Q_T\setminus T|=|M\setminus T|$. By assumption, $A_T>0$, $|M\setminus T|=o(\ln(e^{A_T}))$ and $\mu_{Q_TT}=o(\ln(e^{A_T}))$, then by Lemma~\ref{lemma:s20}, with $c=2$ and $\alpha\in(\psi,1)$, and every $n$ large enough,
\begin{equation}\label{eq:boundLQS}
\int_0^1P\Bigl(L_{Q_TT}>2\ln{\Bigl(\frac{e^{A_T}}{1 / u-1}\Bigl)}\Bigl){\rm d}u \;<\; 
e^{-\alpha A_T}.
\end{equation}

We now bound the second term in~\eqref{eq:prooflemmas20nonspu}. If $M \supset T$, then $Q_T=M$, $L_{Q_TM}=0$, $\mu_{Q_TM}=0$, and this term is zero. Alternatively, if $M \not\supset T$
then, by Lemma~\ref{lemma:s7}, $L_{Q_TM} \sim \chi_{|Q_T\setminus M|}^2(\mu_{Q_TM})$ with $|Q_T\setminus M|=|T\setminus M|$.
Clearly, when $\mu_{Q_TM}\leq 6\Delta_{MT}$ this term is also zero, so suppose that $\mu_{Q_TM}> 6\Delta_{MT}$.  
We have  
$$P\left(L_{Q_TM}<2\gamma'(\mu_{Q_TM}-6\Delta_{MT})\right)\; \leq \; P\left(L_{Q_TM}<2\gamma'\mu_{Q_TM}\right)$$ 
and, by Lemma~\ref{lemma:s2} and using that $\gamma'\in (0,\tfrac{1}{12})$, we obtain that 
$$P\left(L_{Q_TM}<2\gamma'\mu_{Q_TM}\right) \;\;\leq\;\; (2\gamma')^{\frac{|T\setminus M|}{4}} e^{-\frac{1}{2}(1-\sqrt{2\gamma'})^2\mu_{Q_TM}}\;\;\leq\;\;  (\tfrac16)^{\frac{|T\setminus M|}{4}}e^{-\frac{1}{2}(1-\sqrt{2\gamma'})^2\mu_{Q_TM}}.$$
Since $\mu_{Q_TM}>6\Delta_{MT}$, we get
$$
A_T\;=\;\gamma'\mu_{Q_TM}+\gamma\Delta_{MT}\;<\;\frac{1-\gamma}{6}\mu_{Q_TM}+\frac{\gamma}{6}\mu_{Q_TM}\;=\;\frac16\mu_{Q_TM}\;\leq\;\frac12(1-\sqrt{2\gamma'})^2\mu_{Q_TM},
$$
where the last inequality follows from the fact that $\gamma'\in(0, \tfrac{1}{12})$. We then get
\begin{equation}\label{eq:boundLQM}
    P\left(L_{Q_TM}<2\gamma'(\mu_{Q_TM}-6\Delta_{MT})\right) \;\leq\; \left(\frac16\right)^{\frac{|T\setminus M|}{4}}e^{-A_T} \leq e^{-A_T}\leq e^{-\alpha A_T}
\end{equation}
Summing the bounds in~\eqref{eq:boundLQS} and~\eqref{eq:boundLQM} gives that for every $n$ large enough $\E\left(NC(M)\right) < 2e^{-\alpha A_T}$. Since $\psi<\alpha$, for every $n$ large enough, we have that $\E\left(NC(M)\right) < e^{-\psi A_T}$ as claimed.
\end{proof}

\end{document}
